\section{Introduction}
\label{sec:introduction}

An optimizer is an intervention-dependent dynamical system.  Its state may contain parameters, momentum, curvature estimates, preconditioners, schedules, clipping rules, random innovations, target transformations, and numerical state.  A recorded parameter update is one projection of that extended dynamics.  It generally admits many internal explanations: positive geometry may reproduce part of the update, while memory, operators, control, noise, or target changes reproduce the remainder.

This ambiguity creates two distinct scientific questions.  A \emph{response-class question} asks whether an observed trace is close to the outputs generated by a declared family, such as bounded diagonal geometry.  A \emph{mechanism question} asks how the counterfactual response changes when a declared module is altered.  Passive projection can reject a response class.  Causal mechanism effects require interventions, instrumentation, or an observation structure with equivalent identifying power.

We take the complete interventional response as the primitive object.  A lower-finite poset $(\Pcal,\preceq)$ indexes configurations, including Boolean sector masks, graded mechanism levels, and nested geometry budgets.  Under a protocol that fixes initialization, data, horizon, response map, and noise coupling, each configuration $a$ produces a finite-horizon response $F(a)$ in a Hilbert space.  The optimizer may be nonsmooth, stochastic, stateful, delayed, and history-dependent.

Three operations then have separate roles.  Causal realization turns histories into a predictive state.  M\"obius inversion changes coordinates from configuration responses to pure interventional effects.  A linear design operator maps those effects to queried observations.  Their order is typed:
\[
\text{causal dynamics}\longrightarrow
\text{response map}\longrightarrow
\text{M\"obius effects}\longrightarrow
\text{observational quotient}.
\]
The universal theory requires no optimizer energy or geometry.

The companion P1 paper already answers the structural reduced-value question.  When intervention amplitudes couple to hidden optimizer variables that relax by energy minimization, it proves that the reduced intervention Hessian is
\[
D^2_{uu}\bar E=-G^*H^{-1}G,
\]
and that Boolean reduced-value contrasts integrate this curvature \citep{li2026optimizationgeometrodynamics}.  P2 supplies the corresponding path-space M\"obius--Jacobi theorem \citep{li2026restrictedcomplexity}.  The open problem for an experimental calculus is that optimizer studies usually observe updates, traces, or task statistics rather than the reduced optimal value itself.

We solve this interface problem with an observable readout theorem.  For any smooth response $F_R(u)=R(u,\sigma_\star(u))$, its mixed derivative is an explicit five-term transfer through the first and second hidden responses.  The P1 negative Gram law is recovered for the reduced-value readout, while an actual update readout generally has no universal sign.  Its Boolean M\"obius effect remains the exact integral of this transferred curvature and is therefore accessible to a factorial design.  This separates three statements that had previously been conflated: P1 explains reduced-value curvature, the new transfer law predicts the chosen observable, and the incidence experiment determines whether that observable effect is identifiable.

This positions the paper as the experimental outer layer of a four-paper program.  Information-induced training geometry supplies explicit visible metrics and minimal SPD completion \citep{li2026informationinducedgeometry}.  P1 supplies finite-dimensional hidden-state reduction and reduced-value curvature \citep{li2026optimizationgeometrodynamics}.  P2 supplies structured path-space geometry, Jacobi response, and its global M\"obius realization \citep{li2026restrictedcomplexity}.  The present paper contributes the pathwise causal state, observable-readout transfer, experimental gauge, inference, and design layers.  Its universal representation theorem does not require the regular companion theories.

\paragraph{Contributions.}
\begin{enumerate}[leftmargin=*,itemsep=3pt]
\item We prove a universal causal--interventional representation theorem.  Under the protocol's fixed innovation coupling, pathwise predictive equivalence gives the canonical behaviorally minimal optimizer state.  Lower-finite incidence coordinates give unique pure effects, the complete design gauge, exact identifiability, sharp quotient stability, held-out falsification, and exact noiseless configuration complexity for every finite support.
\item Taking the P1 Schur--M\"obius theorem as a structural input, we prove an observable readout interaction law.  It transfers first and second hidden relaxation responses to arbitrary smooth scalar readouts, covers actual update coordinates, integrates exactly to their Boolean effects, and shows why the reduced-value sign law does not extend to general optimizer traces.
\item We develop the identifiable observation layer: projection certificates, Gaussian quotient minimax risk, exact confidence sets, residualized sector tests, misspecification decomposition, certified downstream actions, total-recovery allocation, and exact target-subspace A-optimal replication.  The unrestricted Boolean recovery risk is exponential, quantifying the price of universal attribution.
\item We complete a controlled real-data factorial closure on a 65-dimensional strongly convex logistic hidden problem.  Independent Boolean and curvature paths, continuous held-out intensities, eight data splits, information-optimal allocation, 100,000 Gaussian campaigns, 4500 real-minibatch observations, 20,000 bootstrap draws, held-out support rejection, and finite-response bounds validate the reduced-value structural input and the identification/inference layer.  Separate neural audits instantiate full-SPD, diagonal, block, channel-scalar, layer-scalar, and trajectory-restricted response classes in nonconvex training.
\end{enumerate}

\paragraph{Scope.}
Universality concerns pathwise causal representation under a fixed innovation coupling and incidence identification of expected responses.  Smooth readout interaction requires the stated regularity, uniqueness, and positive vertical-Hessian assumptions.  Exact Gaussian inference additionally assumes a finite-dimensional response and known positive-definite covariance.  Trace explainability is not an optimizer-performance ranking.
